\documentclass[12 pt, letterpaper]{hmcpset}
\usepackage{hw-preamble}
\setlist[enumerate, 1]{label = (\roman*)}

\name{Jayden Thadani}
\class{MATH 145 --- Topics in Geometry and Topology}
\assignment{Homework 3}
\duedate{12th February 2025}

\begin{document}

\begin{problem}[15. (Brouwer's fixed point theorem in 2D)]
	Let $D^{2}$ denote the closed unit disc in the complex plane, and let $f : D^{2} \to D^{2}$ be continuous. Brouwer's theorem asserts that $f$ has at least one \emph{fixed point}, that is, a point $x \in D^{2}$ such that $f(x) = x$.

	Suppose $f$ has no fixed points. Then we can define a continuous function $g : D^{2} \to \mathbb{C} \setminus \{ 0 \}$ by the formula $g(z) = z - f(z)$.
	\begin{enumerate}
		\item Show that the loops
			\[
				t \mapsto g(e^{2 \pi i t}) \quad \text{and} \quad t \mapsto e^{2 \pi i t}
			\]
			are homotopic in $\mathbb{C} \setminus \{ 0 \}$. It follows that $w(g(e^{2 \pi i t}), 0) = 1$.
		\item Show that the loops
			\[
				t \mapsto g(e^{2 \pi i t}) \quad \text{and} \quad t \mapsto g(0)
			\]
			are homotopic in $\mathbb{C} \setminus \{ 0 \}$. It follows that $w(g(e^{2 \pi i t}), 0) = 0$.
	\end{enumerate}
	This contradiction implies the Brouwer fixed point theorem in 2 dimensions.
\end{problem}

\begin{solution}
	We will write $h : t \mapsto e^{2 \pi i t}$.
	\begin{enumerate}
		\item We claim that $g \circ h(z)$ and $h(z)$ always have (the support of) the straight line segment between them contained in $\mathbb{C} \setminus \{ 0 \}$. 

			Since $g \circ h = h - f \circ h$, $\lvert h \rvert = 1$, and $\lvert f \rvert \le 1$, the line segment between  $g \circ h(z)$ and $h(z)$ has length $1$. The only line segment with length less than or equal to $1$ that has one endpoint on the circle $h$ and is not contained in $\mathbb{C} \setminus \{ 0 \}$ is the radius, which has as its other end point $0$. Since $g \circ h(z)$ is never $0$, the straight line between $h$ and $g \circ h$ is always contained in $\mathbb{C} \setminus \{ 0 \}$.

			Thus, there is a straight line homotopy giving $g \circ h \simeq h$ in $\mathbb{C} \setminus \{ 0 \}$.

		\item Notice that $h$ and $0$ are homotopic in $\mathbb{C}$ by straight line homotopy. Since $g$ is a continuous function $\mathbb{C} \to \mathbb{C} \setminus \{ 0 \}$ we have that $g \circ h \simeq g(0)$ in the codomain $\mathbb{C} \setminus \{ 0 \}$.
	\end{enumerate}
\end{solution}

\pagebreak

\begin{problem}[16.]
	Let's define a \emph{Brouwer space} to be a topological space $X$ such that every continuous map $f : X \to X$ has at least one fixed point. It is a known theorem that the closed unit disc in $n$-dimensional space, $D^{n}$ is a Brouwer space for every $n$.

	Let us now study some properties of Brouwer spaces.
	\begin{enumerate}
		\item Let $Y$ be a Brouwer space. Suppose $X$ is homeomorphic to $Y$. This means that there are continuous functions
			\[
				p : X \to Y \quad \text{and} \quad q : Y \to X
			\]
			such that
			\[
				q \circ p = \operatorname{id}_{X} \quad \text{and} \quad p \circ q = \operatorname{id}_{Y}.
			\]
			Show that $X$ is a Brouwer space.
		\item Let $Y$ be a Brouwer space. Suppose $X$ is a \emph{retract} of $Y$. This means that $X$ is a subspace of $Y$, with inclusion map written $p : X \to Y$ and there is a continuous function $q : Y \to X$ such that $q \circ p = \operatorname{id}_{X}$.

			Show that $X$ is a Brouwer space.
	\end{enumerate}
\end{problem}

\begin{solution}
	\begin{enumerate}
		\item Notice that if $X$ is homeomorphic to $Y$, then it is a retract of $Y$ since $p$ is still an inclusion (an injective continuous map) and $q \circ p = \operatorname{id}_{X}$. Thus, we can use part (ii) to show that this is true.
		\item Suppose $f : X \to X$ is continuous. Notice that $g = p \circ f \circ q$ is continuous too and the composition $Y \to X \to X \to Y$ yields a map $Y \to Y$. That is, the diagram below commutes
			\[
				\begin{tikzcd}
					X \ar[r, "f"] & X \ar[d, "p"] \\
					Y \ar[u, "q"] \ar[r, "g"] & Y
				\end{tikzcd}
			\]

			Since  $Y$ is a Brouwer space, $g$ has a fixed point $y \in Y$. We claim that $q(y)$ is a fixed point of $f$ --- 

			We can see this
			\[
				f(q(y)) = (\operatorname{id}_{X} \circ f \circ q)(y) = q \circ (p \circ f \circ q)(y) = q(g(y)) = q(y)
			\]
			where the second equality is because $\operatorname{id}_{X} = q \circ p$ and the third equality is because we defined $g = p \circ f \circ q$.

			We could alternatively argue that $q(y) = y$ since $g$ has image contained in $X$, and thus, any fixed point in $Y$ must be a fixed point in $X$.
	\end{enumerate}
\end{solution}

\pagebreak

\begin{problem}[17.]
	\begin{enumerate}
		\item The unit sphere in $n$-dimensions
			\[
				S^{n - 1} = \{ x \in \mathbb{R}^{n} \mid \lVert x \rVert = 1 \}
			\]
			is a subspace of the annulus
			\[
				A^{n} = \{ x \in \mathbb{R}^{n} \mid 1/2 \le \lVert x \rVert \le 1 \}.
			\]
			Show that $S^{n - 1}$ is a retract of $A^{n}$.
		\item Show that $S^{n - 1}$ is not a Brouwer space (for any $n$)
		\item The unit sphere $S^{n - 1}$ is a subspace of the closed unit disc
			\[
				D^{n} = \{ x \in \mathbb{R}^{n} \mid \lVert x \rVert \le 1 \}.
			\]

			Assuming the general Brouwer theorem, show that $S^{n - 1}$ is not a retract of $D^{n}$.
	\end{enumerate}
\end{problem}

\begin{solution}
	\begin{enumerate}
		\item Since $S^{n - 1}$ is a subspace of the annulus, we have the inclusion $p : S^{n - 1} \to A^{n}$. For $q$ we dilate each point in the annulus onto the sphere. That is,
			\[
				q : x \mapsto \frac{x}{\lVert x \rVert}.
			\]
			This is the continuous quotient of the identity function and the norm and thus, continuous. Since all $x \in S^{n - 1}$ have $\lVert x \rVert = 1$,  $q \circ p = \operatorname{id}_{S^{n - 1}}$.

		\item Consider the continuous map $x \mapsto -x$ on $\mathbb{R}^{n}$ restricted to $S^{n - 1}$. Since $\lVert x \rVert = \lVert -x \rVert$ this is a map $S^{n - 1} \to S^{n - 1}$. However, $0 \notin S^{n - 1}$ which is the only fixed point of the map on $\mathbb{R}^{n}$. Thus, the map does not have a fixed point --- $S^{n - 1}$ is not a Brouwer space.

		\item By Brouwer's fixed point theorem $D^{n}$ is a Brouwer space. If $S^{n - 1}$ were a retract of $D^{n}$ we know by problem 16 that it must also be a Brouwer space. In the previous subpart we have shown that $S^{n - 1}$ is not a Brouwer space. Thus, $S^{n - 1}$ cannot be a retract of $D^{n}$.
	\end{enumerate}
\end{solution}

\pagebreak

\begin{problem}[18.]
	\begin{enumerate}
		\item Let $Q$ denote the closed square disc
			\[
				\{ (x, y) \in \mathbb{R}^{2} \mid \max \{ \lvert x \rvert, \lvert y \rvert \} \le 1 \}
			\]
			and let $X$ denote its cross-shaped subspace
			\[
				X = \{ (x, y) \in Q \mid \lvert x \rvert = \lvert y \rvert \}
			\]

			Show that $X$ is a retract of $Q$, by writing down a suitable continuous map $Q \to X$.

		\item It follows from (i) that the cross-shaped space $X$ is a Brouwer space. Sketch or otherwise describe an entertainingly mildly complicated continuous function $f : X \to X$ and indicate at least one fixed point.
	\end{enumerate}
\end{problem}

\begin{solution}
	\begin{enumerate}
		\item Notice that the hypotenuse of a right-angled triangle is a retract of the whole triangle. For concreteness, consider the triangle $\Delta = \{ (x, y) \in \mathbb{R}^{2} \mid x > 0, y > 0, x + y \le 1 \}$. The continuous map $(x, y) \mapsto (x, 1 - x)$ is the retraction map. By dividing the polydisc into eight triangles and orienting them so that the continuous on each triangle on intersections, we get a retraction map $q : Q \to X$.

			We can give a direct formula for $q$ as well.
			\[
				q(x, y) = \begin{cases}
					(x, 1 - \lvert x \rvert) & y > 0, y > \lvert x \rvert \\
					(x, -1 + \lvert x \rvert) & y < 0, \lvert y \rvert > \lvert x \rvert \\
					(1 - \lvert y \rvert, y) & x > 0, x > \lvert y \rvert \\
					(-1 + \lvert y \rvert, y) & x < 0, \lvert x \rvert > \lvert y \rvert.
				\end{cases}
			\]

			\begin{figure}[H]
				\centering
				\includegraphics[width = 0.4 \textwidth]{figures/P18i retract}
				\caption{Arrows indicating the flow of the retraction map $q : Q \to X$}
			\end{figure}

		\item We can try to avoid fixed points except $0$ by first flipping both arms of $X$ --- mapping $(x, y) \mapsto -(x, y)$. Then to remove the fixed point at $0$ we can compose the flipping with squishing one of the arms towards one end (given a homeomorphism $f : [0, 1] \to \{ (x, y) \in Q \mid x = y \}$ with $f(1/2) = 0$ consider $f \circ x^{2} \circ f^{-1}$) and the other arm accordingly so that the function is continuous. However, if we squish all of $X$ in one direction there will be a point that was close to $(0, 0)$ before flipping (on the side towards which we squish the arm) and gets flipped away and then squished back to itself under squishing. In our example of squishing, $0.5 - 0.12 \approx (0.5 + 0.12)^{2}$.
	\end{enumerate}
\end{solution}

\pagebreak

\begin{problem}[19.]
	\begin{enumerate}
		\item Explain how to interpret the following diagram as a $1$-cycle with coefficients in $\mathbb{F}_{2}$.

			Specifically, label the vertices $a, b, c, d, e$ in counter-clockwise order from the top. Write down the $1$-cycle in the form  $\sum \lambda_{i} [a_{i}, b_{i}]$ and confirm that it is a $1$-cycle in $\mathbb{F}_{2}$.
		\item Calculate the $\mathbb{F}_{2}$ winding number $w_{\mathbb{F}_{2}}$ of your cycle in each region.
	\end{enumerate}
\end{problem}

\begin{solution}
	\begin{enumerate}
		\item We choose the order-preserving labeling --- if $a_{1} < a_{2}$ lexicographically, we orient the edge between them $[a_{1}, a_{2}]$. Thus, our $1$-cycle is
			\[
				[a, b] + [a, c] + [a, d] + [a, e] + [b, c] + [b, d] + [b, e] + [c, d] + [c, e] + [d, e]
			\]
			(all the coefficients are $1$). This works because $K_{5}$ has an odd number of vertices and thus, each vertex has even degree. If $I$ is the number of edges pointing in and $O$ is the number pointing out, we must have $I + O = 4$, and thus, $I - O = I + O - 2 O$ is even too. That is, $I + O$ is sent to $0$ in the canonical map $\mathbb{Z} \to \mathbb{F}_{2}$.

			Notice that this proof works for any orientation of $K_{5}$ that has nonzero coefficient for all of its edges. Another way to see this is that reversing orientation is equivalent to negating coefficients since $-[a, b] = [b, a]$, but $1 = -1$ in $\mathbb{F}_{2}$.

		\item The winding numbers below were calculated using the ray escape formula. Since $-1 = 1$ in $\mathbb{F}_{2}$, the sign of the crossing doesn't matter --- the winding number is determined by the parity of the number of lines crossed. 

			The outermost triangles require one crossing, the inner triangles require two, and the pentagon requires three. The winding number in the unbounded region is $0$ by the distant cycle lemma. Thus, the winding numbers are as follows ---
			\begin{figure}[H]
				\centering
				\includegraphics[width = 0.4 \textwidth]{figures/P19i K5}
				\caption{The winding numbers in each (bounded) connected region of the complement $K_{5}$ are written in blue}
			\end{figure}
	\end{enumerate}
\end{solution}

\pagebreak

\begin{problem}[20.]
	\begin{enumerate}
		\item Show that the twelve edges of the following graph can be oriented to form a $1$-cycle over the field $\mathbb{F}_{3} = \mathbb{Z} / 3 \mathbb{Z}$.
		\item Calculate the $\mathbb{F}_{3}$-winding number $w_{\mathbb{F}_{3}}$ of your cycle in each region.
	\end{enumerate}
\end{problem}

\begin{solution}
	One can verify that each vertex in the graph below either has outdegree $3$ or indegree $3$. Both are sent to zero $0$ in the canonical map $\mathbb{Z} \to \mathbb{F}_{3}$.

	The winding numbers were calculated using the ray escape formula again. Note that in $\mathbb{F}_{3}$ we have $-1 = 2$ so regions with clockwise crossings have winding number $2$.
			\begin{figure}[H]
				\centering
				\includegraphics[width = 0.5 \textwidth]{figures/P20i cube.pdf}
				\caption{The winding numbers in each (bounded) connected region of the complement of the cube are written in blue}
			\end{figure}
\end{solution}

\pagebreak

\begin{problem}[21.]
	We have three loops
	\[
		f, g, h : [0, 1] \to \mathbb{R}^{2} \setminus \{ 0 \}
	\]
	that are related in the following way.

	Let $D$ denote this region in the plane:

	\begin{figure}[H]
		\centering
		\includegraphics[width = 0.4 \textwidth]{figures/P21 D8.pdf}
	\end{figure}

	Let $\alpha, \beta, \psi : [0, 1] \to D$ be loops defined as follows
	\begin{itemize}
		\item $\alpha$ traverses the small interior boundary square on the left side of $D$ counter-clockwise at uniform speed, starting and ending at the square's bottom right corner.
		\item $\beta$ traverses the small interior boundary square on the right hand side of $D$, counter-clockwise at uniform speed, starting and ending at the square's bottom right corner.
		\item $\gamma$ traverses the rectangular outer boundary of $D$, counter-clockwise at uniform speed, starting and ending at the rectangle's bottom right corner.
	\end{itemize}

	Suppose we have a map $H : D \to \mathbb{R}^{2} \setminus 0$ such that
	\[
		f(t) = H(\alpha(t)), \quad g(t) = H(\beta(t)), \quad h(t) = H(\gamma(t))
	\]
	for all $t \in [0, 1]$. Show that the winding numbers are related by the formula
	\[
		w(f, 0) + w(g, 0) = w(h, 0)
	\]
	by considering a subdivision of $D$ into $13 N^{2}$ equal squares where $N$ is sufficiently large.
\end{problem}

\begin{solution}
	Let $m$ be the minimum distance of $H$ from $0$ --- that is, $m = \min \{ \lVert H(\mathbf{x}) \rVert \mid \mathbf{x} \in D \}$. Note that this minimum exists and is positive since each of the compositions with $H$ is a continuous function from the compact set $D$ to $\mathbb{R}^{2} \setminus \{ 0 \}$.

	Since $H$ is a (uniformly) continuous function on a compact set, we can choose $\delta$ such that for all $\mathbf{x}_{1}, \mathbf{x}_{2} \in D$ with $\lVert \mathbf{x}_{2} - \mathbf{x}_{1} \rVert < \delta$, we have
	\[
		\lVert H(\mathbf{x}_{2}) - H(\mathbf{x}_{1}) \rVert < m / 2
	\]

	We choose $N$ so that $1 / N < \delta$, divide $D$ into thirteen unit squares, and further subdivide each of those squares into the $N^{2}$ squares, each of side length $1 / N$ like so ---
	\begin{figure}[H]
		\centering
		\includegraphics[width = 0.5 \textwidth]{figures/P21 D8 annotated}
		\caption{The figure 8 $D$ divided into $13$ squares and each square further subdivided into $N^{2}$ parts, with curves $\alpha$, $\beta$, and $\gamma$}
	\end{figure}

	Now we can replace each of $\alpha, \beta, \gamma$ with the $1$-cycles $\alpha_{T}, \beta_{T}, \gamma_{T}$ obtained by sampling their points on the partition of $D$. Since the partition we chose is fine enough, the winding numbers of $H \circ \alpha_{T}, H \circ \beta_{T}, H \circ \gamma_{T}$ around $0$ don't change by refining any further.

	Consider all the squares (positively oriented) between $\gamma_{T}$ and $\alpha_{T}, \beta_{T}$. We claim that under mapping by $H$, their total winding number around $0$ is $0$, and thus, the winding number of $H \circ \gamma_{T}$ is just the sum of the winding numbers of $H \circ \alpha_{T}$ and $H \circ \beta_{T}$. 

	This follows from the distant cycle lemma. Given any square with vertices $\mathbf{a}, \mathbf{b}, \mathbf{c}, \mathbf{d}$ we have $H(\mathbf{a}) \ge m$ and by the triangle inequality (this is why we used $m / 2$ to define $\delta$ earlier) the distance between each pair of $H(\mathbf{a}), H(\mathbf{b}), H(\mathbf{c}), H(\mathbf{d})$ is less than $m$. Thus, by the distant cycle lemma, we have that the cycle of the square (pushed forward by $H$ into $\mathbb{R}^{2}$) has winding number $0$ around $0$.

	By cancelling out the common edges of the squares between $\gamma_{T}$ and $\alpha_{T}, \beta_{T}$, since we oriented the squares counterclockwise, we're left with a $1$-chain $\Gamma = \gamma_{T} - (\alpha_{T} + \beta_{T})$ and $w(H \circ \Gamma, 0) = 0$. By additivity, $w(H \circ \gamma_{T}) = w(H \circ \alpha_{T}) + w(H \circ \beta_{T})$, and since the samplings on the partition have the same winding number as the actual curves, we have
	\[
		w(H \circ \gamma, 0) = w(H \circ \alpha, 0) + w(H \circ \beta, 0)
	\]
	which is just
	\[
		w(f, 0) + w(g, 0) = w(h, 0)
	\]
	as desired.
\end{solution}

\end{document}
